% Options for packages loaded elsewhere
\PassOptionsToPackage{unicode}{hyperref}
\PassOptionsToPackage{hyphens}{url}
\PassOptionsToPackage{dvipsnames,svgnames,x11names}{xcolor}
\documentclass[
  10pt,
  fontset=fandol]{ctexart}
\usepackage{xcolor}
\usepackage[margin=16mm]{geometry}
\usepackage{amsmath,amssymb}
\setcounter{secnumdepth}{-\maxdimen} % remove section numbering
\usepackage{iftex}
\ifPDFTeX
  \usepackage[T1]{fontenc}
  \usepackage[utf8]{inputenc}
  \usepackage{textcomp} % provide euro and other symbols
\else % if luatex or xetex
  \usepackage{unicode-math} % this also loads fontspec
  \defaultfontfeatures{Scale=MatchLowercase}
  \defaultfontfeatures[\rmfamily]{Ligatures=TeX,Scale=1}
\fi
\usepackage{lmodern}
\ifPDFTeX\else
  % xetex/luatex font selection
\fi
% Use upquote if available, for straight quotes in verbatim environments
\IfFileExists{upquote.sty}{\usepackage{upquote}}{}
\IfFileExists{microtype.sty}{% use microtype if available
  \usepackage[]{microtype}
  \UseMicrotypeSet[protrusion]{basicmath} % disable protrusion for tt fonts
}{}
\makeatletter
\@ifundefined{KOMAClassName}{% if non-KOMA class
  \IfFileExists{parskip.sty}{%
    \usepackage{parskip}
  }{% else
    \setlength{\parindent}{0pt}
    \setlength{\parskip}{6pt plus 2pt minus 1pt}}
}{% if KOMA class
  \KOMAoptions{parskip=half}}
\makeatother
\setlength{\emergencystretch}{3em} % prevent overfull lines
\providecommand{\tightlist}{%
  \setlength{\itemsep}{0pt}\setlength{\parskip}{0pt}}
\usepackage{amsmath,amssymb,mathtools,needspace}
\xeCJKDeclareCharClass{CJK}{"2032,"0400->"04FF,"2160->"216B,"2460->"2473,"25A0,"25C7,"25CB,"25CF}
\IfFontExistsTF{Arial Unicode MS}{\xeCJKsetup{AutoFallBack=true}\setCJKfallbackfamilyfont{\CJKrmdefault}{Arial Unicode MS}}{}
\setcounter{secnumdepth}{0}
\setlength{\emergencystretch}{2em}
\allowdisplaybreaks
\usepackage{bookmark}
\IfFileExists{xurl.sty}{\usepackage{xurl}}{} % add URL line breaks if available
\urlstyle{same}
\hypersetup{
  pdftitle={基于 Taylor 展开的构造方法},
  colorlinks=true,
  linkcolor={Maroon},
  filecolor={Maroon},
  citecolor={Blue},
  urlcolor={Blue},
  pdfcreator={LaTeX via pandoc}}

\title{基于 Taylor 展开的构造方法}
\author{}
\date{}

\begin{document}
\maketitle

\subsubsection{5.5.5 基于 Taylor 展开的构造方法}\label{ux57faux4e8e-taylor-ux5c55ux5f00ux7684ux6784ux9020ux65b9ux6cd5}

我们看到，基于数值积分可以构造出一系列求解常微分方程的计算格式，然而下面将要介绍的 Taylor 展开方法更具有一般性，各种线性多步格式（包括前述 Adams 格式）均可通过这种方法构造出来。

一般的线性多步格式具有下列形式：

\[
y_{n+1}=\alpha_0y_n+\alpha_1y_{n-1}+\cdots+\alpha_ry_{n-r}+h(\beta_{-1}y'_{n+1}+\beta_0y'_n+\beta_1y'_{n-1}+\cdots+\beta_ry'_{n-r}),
\]

缩记作

\begin{quote}
校注（agent 补充，ch05b-E002）：本页两条``事后估计式''及计算方案两条``改进''式的分母均印为 \(720\)，已照录。由本页前两条局部误差式相减，\(y_{n+1}-\bar y_{n+1}\approx(270/720)h^5y^{(5)}(x_n)\)，消去 \(h^5y^{(5)}(x_n)\) 后，事后估计的系数应为 \(251/270\) 和 \(19/270\)；相应改进式中的分母也应为 \(270\)。这是原书疑误，非 OCR 错。

校注（agent 补充，ch05b-E003）：计算方案第一条``计算''式左端印为 \(m_{n+1}\)，未带撇号；已照录。按下一条校正式使用的 \(m'_{n+1}\) 及导数定义，此处疑漏印撇号，应为 \(m'_{n+1}=f(x_{n+1},m_{n+1})\)。

校注（agent 补充，ch05b-E004）：表 5.10 末行准确值原印为 \(1.7231\)，已照录。已另行目视核对 PDF 120（书 107 页）式（5.2.2）后所列解析解 \(y=\sqrt{1+2x}\) 及表 5.1；\(y(1)=\sqrt3\approx1.7321\)，因此本格疑有数字对调。
\end{quote}

\href{../../source-images/pdf-142.jpeg}{原页}

\[
y_{n+1}=\sum_{k=0}^r\alpha_ky_{n-k}+h\sum_{k=-1}^r\beta_ky'_{n-k},
\tag{5.5.14}
\]

其中 \(y'_{n-k}=f(x_{n-k},y_{n-k})\)。注意到 \(y'_{n+1}=f(x_{n+1},y_{n+1})\) 中含有未知的 \(y_{n+1}\)，则当 \(\beta_{-1}=0\) 时，格式（5.5.14）是显式的，\(\beta_{-1}\ne0\) 时则是隐式的。

设 \(y_{n-k}=y(x_{n-k}),y'_{n-k}=y'(x_{n-k})\)，由 Taylor 展开，有

\[
y_{n-k}=\sum_{j=0}^p\frac{(-kh)^j}{j!}y_n^{(j)}+\frac{(-kh)^{p+1}}{(p+1)!}y_n^{(p+1)}+\cdots,
\]

\[
y'_{n-k}=\sum_{j=1}^p\frac{(-kh)^{j-1}}{(j-1)!}y_n^{(j)}+\frac{(-kh)^p}{p!}y_n^{(p+1)}+\cdots.
\]

代入式（5.5.14），整理得

\[
\begin{aligned}
y_{n+1}={}&\left(\sum_{k=0}^r\alpha_k\right)y_n+\sum_{j=1}^p\frac{h^j}{j!}\left[\sum_{k=1}^r(-k)^j\alpha_k+j\sum_{k=-1}^r(-k)^{j-1}\beta_k\right]y_n^{(j)}\\
&+\frac{h^{p+1}}{(p+1)!}\left[\sum_{k=1}^r(-k)^{p+1}\alpha_k+(p+1)\sum_{k=-1}^r(-k)^p\beta_k\right]y_n^{(p+1)}+\cdots.
\end{aligned}
\tag{5.5.15}
\]

于是，欲使格式（5.5.14）成为 \(p\) 阶精度的，即局部截断误差为 \(O(h^{p+1})\)，只要令展开式（5.5.15）与 \(y(x_{n+1})\) 的 Taylor 展开式

\[
y(x_{n+1})=\sum_{j=0}^p\frac{h^j}{j!}y_n^{(j)}+\frac{h^{p+1}}{(p+1)!}y_n^{(p+1)}+\cdots
\tag{5.5.16}
\]

能符合到 \(h^p\) 项，为此要求成立

\[
\begin{cases}
\displaystyle\sum_{k=0}^r\alpha_k=1,\\[6pt]
\displaystyle\sum_{k=1}^r(-k)^j\alpha_k+j\sum_{k=-1}^r(-k)^{j-1}\beta_k=1\qquad(j=1,2,\cdots,p).
\end{cases}
\tag{5.5.17}
\]

如果格式（5.5.14）的系数满足这组条件，则将式（5.5.16）与式（5.5.15）相减，即得局部截断误差

\[
\begin{aligned}
&y(x_{n+1})-y_{n+1}\\
&=\frac{h^{p+1}}{(p+1)!}\left[1-\sum_{k=1}^r(-k)^{p+1}\alpha_k-(p+1)\sum_{k=-1}^r(-k)^p\beta_k\right]y_n^{(p+1)}+\cdots.
\end{aligned}
\tag{5.5.18}
\]

下面进一步具体地考察四步方法。先研究下列形式的四步显式格式：

\[
y_{n+1}=\alpha_0y_n+\alpha_1y_{n-1}+\alpha_2y_{n-2}+h(\beta_0y'_n+\beta_1y'_{n-1}+\beta_2y'_{n-2}+\beta_3y'_{n-3}),
\tag{5.5.19}
\]

欲使这类格式成为四阶的，按式（5.5.17），其系数应当满足条件

\[
\begin{cases}
\alpha_0+\alpha_1+\alpha_2=1,\\
-\alpha_1-2\alpha_2+\beta_0+\beta_1+\beta_2+\beta_3=1,\\
\alpha_1+4\alpha_2-2\beta_1-4\beta_2-6\beta_3=1,\\
-\alpha_1-8\alpha_2+3\beta_1+12\beta_2+27\beta_3=1,\\
\alpha_1+16\alpha_2-4\beta_1-32\beta_2-108\beta_3=1.
\end{cases}
\tag{5.5.20}
\]

\href{../../source-images/pdf-143.jpeg}{原页}

这里有七个待定系数 \(\alpha_0,\allowbreak \alpha_1,\allowbreak \alpha_2,\allowbreak \beta_0,\allowbreak \beta_1,\allowbreak \beta_2,\allowbreak \beta_3\)，然而它们所要适合的条件只有五个，因此有两个自由度。设令 \(\alpha_1=\alpha_2=0\)，求解式（5.5.20）得到

\[
\alpha_0=1,\quad\beta_0=\frac{55}{24},\quad\beta_1=-\frac{59}{24},\quad\beta_2=\frac{37}{24},\quad\beta_3=-\frac9{24}.
\]

这时式（5.5.19）为四阶 Adams 格式（5.5.6）。另外，用定出的系数值具体地代入式（5.5.18），即可导出误差公式（5.5.10）。

为了提供与式（5.5.19）相匹配的隐式方法，舍弃其中的 \(y'_{n-3}\) 而代之以 \(y'_{n+1}\)，有

\[
y_{n+1}=\alpha_0y_{n+1}\alpha_1y_{n-1}+\alpha_2y_{n-2}+h(\beta_{-1}y'_{n+1}+\beta_0y'_n+\beta_1y'_{n-1}+\beta_2y'_{n-2}).
\tag{5.5.21}
\]

为使这类格式具有四阶精度，按式（5.5.17），系数应当满足条件

\[
\begin{cases}
\alpha_0+\alpha_1+\alpha_2=1,\\
-\alpha_1-2\alpha_2+\beta_{-1}+\beta_0+\beta_1+\beta_2=1,\\
\alpha_1+4\alpha_2+2\beta_{-1}-2\beta_1-4\beta_2=1,\\
-\alpha_1-8\alpha_2+3\beta_{-1}+3\beta_1+12\beta_2=1,\\
\alpha_1+16\alpha_2+4\beta_{-1}-4\beta_1-32\beta_2=1.
\end{cases}
\tag{5.5.22}
\]

特别地，取 \(\alpha_1=\alpha_2=0\) 即得四阶 Adams 格式（5.5.9），通过这一途径同时可以导出误差公式（5.5.11）。

\begin{center}\rule{0.5\linewidth}{0.5pt}\end{center}

\subsection{相关原页校注（agent 补充）}\label{ux76f8ux5173ux539fux9875ux6821ux6ce8agent-ux8865ux5145}

以下校注来自本节覆盖的原页，可能也涉及同页相邻小节；原文照录与校正说明分开保留。

\subsubsection{原书 PDF 页 143 的校注}\label{ux539fux4e66-pdf-ux9875-143-ux7684ux6821ux6ce8}

\href{../pages/pdf-143.md}{完整原页转录} · \href{../../source-images/pdf-143.jpeg}{原页图像}

校注记录：ch05b-E005。

\begin{quote}
校注（agent 补充，ch05b-E005）：式（5.5.21）右端开头在原图中印作 \(\alpha_0y_{n+1}\alpha_1y_{n-1}\)（``\(+1\)''下沉至下标位置，且两项之间无正常加号），上式按可见字形保留。由本页``舍弃 \(y'_{n-3}\) 而代之以 \(y'_{n+1}\)''及式（5.5.19）、（5.5.22），此处疑为排印错误，预期应是 \(\alpha_0y_n+\alpha_1y_{n-1}\)。局部原图如下，仅用于保存这一排印疑误的可审计证据。
\end{quote}

\subsection{本节覆盖原页的校注记录}\label{ux672cux8282ux8986ux76d6ux539fux9875ux7684ux6821ux6ce8ux8bb0ux5f55}

下列记录按原页关联，含同页相邻小节的校注；计算时同时读取上文校注和完整原页。

\begin{itemize}
\tightlist
\item
  \texttt{ch05b-E002}：原书 PDF 页 141
\item
  \texttt{ch05b-E003}：原书 PDF 页 141
\item
  \texttt{ch05b-E004}：原书 PDF 页 141
\item
  \texttt{ch05b-E005}：原书 PDF 页 143
\end{itemize}

\end{document}
